% Figura: dos formas de generar un reporte con un VLM, de extremo a extremo o condicionado en evidencia
% externa. Se usa en sec:marco_rep_generacion.

\begin{figure}[H]
    \centering
    \begin{tikzpicture}[x=1mm, y=1mm]

        % (a) De extremo a extremo.
        \node[font=\small\bfseries, anchor=west] at (-6, 30) {(a) De extremo a extremo};
        \node[figentrada, text width=22mm] (a_img) at (8, 18) {Imagen e\\instrucción};
        \node[figtarea, text width=26mm] (a_vlm) at (62, 18) {\gls{vlm}};
        \node[figsalida, text width=26mm] (a_rep) at (122, 18) {Reporte};
        \draw[figflecha] (a_img) -- (a_vlm);
        \draw[figflecha] (a_vlm) -- (a_rep);

        % (b) Condicionado en evidencia externa. Arriba, el componente que aporta la evidencia y la regla de
        % decisión; abajo, el mismo recorrido de (a) con la evidencia entrando al VLM desde arriba.
        \node[font=\small\bfseries, anchor=west] at (-6, 4) {(b) Condicionado en evidencia externa};
        \node[figtarea, text width=28mm] (b_ev) at (38, -8) {Clasificador o\\recuperación};
        \node[figtarea, text width=34mm] (b_pol) at (84, -8)
            {Regla de decisión\\[1pt]{\footnotesize afirmar, diferencial\\o derivar}};
        \node[figentrada, text width=22mm] (b_img) at (8, -34) {Imagen};
        \node[figtarea, text width=26mm] (b_vlm) at (84, -34) {\gls{vlm}};
        \node[figsalida, text width=26mm] (b_rep) at (122, -34) {Reporte};
        \draw[figflecha] (b_img.north) |- (b_ev.west);
        \draw[figflecha] (b_ev) -- (b_pol);
        \draw[figflecha] (b_pol) -- node[figetiqueta, right=1mm] {evidencia $c(x)$} (b_vlm);
        \draw[figflecha] (b_img) -- (b_vlm);
        \draw[figflecha] (b_vlm) -- (b_rep);

    \end{tikzpicture}
    \caption{Dos formas de generar un reporte: el \gls{vlm} solo o con la evidencia de
    otro componente.}
    \label{fig:marco_generacion_reportes}
\end{figure}
